\chapter{Diskussion und Ausblick}
\label{ch:DiskussionAusblick}


\section{Interpretation der Ergebnisse}
\label{sec:InterpretationErgebnisse}


\section{Limitationen und Herausforderungen}
\label{sec:LimitationenHerausforderungen}


\section{Mögliche Weiterentwicklungen / Potenzielle Erweiterungen / Weiterführende Konzepte}
\label{sec:Konzepte}
Im Zuge des Projekts wurden drei verschiedene Konzepte erarbeitet, die darauf abzielen, den Generator durch 
zusätzliche Trainingsstrategien zu verbessern und dadurch die Qualität der erzeugten Daten für nachgelagerte 
Verfahren zu erhöhen. Die folgenden Abschnitte stellen diese Konzepte mit ihren jeweiligen Grundideen sowie den 
identifizierten Vor- und Nachteilen vor. Nach eingehender Diskussion würde die Projektgruppe, sollte das Thema in 
zukünftigen Arbeiten weiterverfolgt werden, insbesondere Konzept~3 als vielversprechendste Option betrachten.

\subsection{Konzept 1: Feature-Anzahl}
\label{sec:Konzept1}
Das erste Konzept zur Kopplung des Generators mit XFeat basiert auf der Annahme, dass eine höhere Anzahl an erkannten Merkmalspunkten ein Indikator für
eine bessere Bildrepräsentation ist. Der Generator aus dem ICCV-Repository erzeugt hierzu synthetisch augmentierte Bilddaten, welche anschließend durch
XFeat verarbeitet werden. Die dabei entstehenden Keypoint-Heatmaps dienen als Signalquelle für ein zusätzliches Lernziel. Anstatt lediglich die nach
Non-Maximum Suppression ausgewählten Keypoints zu zählen, wird auf die vollständige, kontinuierliche Heatmap zurückgegriffen. Durch eine soft-differenzierbare
Approximation der Feature-Anzahl, beispielsweise über eine Sigmoid-Gewichtung oberhalb eines adaptiven Schwellwerts, kann ein „soft count“ berechnet werden.
Dieser Wert wird in Relation zu einer gewünschten Zielanzahl an Merkmalen gesetzt und als zusätzlicher Loss-Term in das Training des Generators integriert.
Der Generator erhält somit eine Belohnung, wenn die Anzahl der detektierten Features oberhalb eines definierten Schwellwerts liegt, und eine Bestrafung, wenn
diese darunterliegt.

Die Idee, Informationen aus Keypoint-Heatmaps zur Steuerung des Trainings zu nutzen, ist als Ansatz etabliert. Arbeiten wie UKPGAN~\cite{you2022ukpgan}
verfolgen einen ähnlichen Ansatz, indem sie generative Modelle zur Kontrolle der Spärlichkeit und Verteilung von Keypoints einsetzen. Zudem existieren
Reinforcement-Learning-gestützte Ansätze wie RIPE~\cite{kunzel2025ripe}, bei denen Reward-Signale ohne explizite Ground-Truth-Daten genutzt werden, um die
Anzahl und Qualität der detektierten Punkte zu steuern.

In Bezug auf die Zielsetzung dieses Projekts – die zuverlässige Bestimmung von Greifpositionen und Ausrichtungen für Roboter an überwiegend texturarmen
Bauteilen – weist das Konzept allerdings sowohl Vorteile als auch Einschränkungen auf. Der zentrale Vorteil liegt in der einfachen Implementierbarkeit:
Mit nur geringen Anpassungen im Trainingsloop des Generators lässt sich das zusätzliche Lernsignal integrieren. Zudem bleibt XFeat vollständig eingefroren,
sodass keine tiefgreifenden Änderungen an der Feature-Extraktion erforderlich sind. Kritisch ist jedoch, dass die reine Anzahl an Keypoints kein
Qualitätskriterium darstellt. Ein Generator könnte versuchen, die Feature-Anzahl durch die Erzeugung künstlicher Texturartefakte zu erhöhen („Reward Hacking“),
was für das geplante Einsatzszenario kontraproduktiv wäre. Um dieses Risiko zu minimieren, ist es notwendig, regulierende Zusatzverluste einzusetzen.

Insgesamt bietet Konzept 1 eine solide Ausgangsbasis, um die Machbarkeit einer differentiellen Kopplung zwischen Generator und Detektor zu demonstrieren.
Für das konkrete Ziel – die Fokussierung auf kantenbasierte Merkmale – ist es jedoch nicht hinreichend und wird deshalb nicht weiter verfolgt.

\subsection{Konzept 2: Loss}
\label{sec:Konzept2}
Die Grundidee besteht darin, Informationen aus der Keypoint-Detektion und -Beschreibung zu nutzen, um den Generator 
gezielt zu steuern. Dabei wird ein geringer XFeat-bezogener Loss als Belohnung verstanden, während ein hoher Loss 
eine Bestrafung darstellt. Auf diese Weise soll der Generator lernen, deformierte Bilder so zu erzeugen, dass darin 
robuste und konsistente Features erkennbar sind, die sowohl für Matching-Aufgaben als auch für visuelle 
Servoing-Strategien nutzbar bleiben~\cite{potje2024xfeat}.

Für die praktische Umsetzung lässt sich dieses Konzept nur dann realisieren, wenn der Generator eng mit XFeat 
gekoppelt wird. Dies bedeutet, dass die Generatorausgabe als Eingabe in XFeat genutzt wird und die daraus berechneten
Keypoints oder Feature-Karten als Grundlage für den Loss dienen. XFeat selbst bleibt dabei eingefroren und wird nicht 
trainiert, sondern liefert ausschließlich das Feedbacksignal. Der Gradientenfluss erfolgt ausschließlich durch den 
Generator, während XFeat als „Sensor“ fungiert, der die Qualität der generierten Strukturen misst. Diese Kopplung 
kann in zwei Varianten umgesetzt werden.

In einer ersten Variante wird XFeat als Black-Box betrachtet. Die Berechnung der Keypoints und Matches erfolgt ohne 
Gradientenfluss und liefert lediglich die Positionen im Bild, an denen differenzierbare Ersatzgrößen wie Kantenstärke, 
Gradientenrichtung oder Patch-Ähnlichkeiten bewertet werden. Diese Größen werden zu einem Score zwischen 0 und 1 
kombiniert, der angibt, wie konsistent die Merkmale zwischen Original- und deformiertem Bild sind. Der resultierende 
Loss kann anschließend direkt in die Gesamtverlustfunktion des Generators integriert werden. Der Vorteil dieser Variante 
liegt in ihrer Robustheit und einfachen Umsetzbarkeit, während der Nachteil darin besteht, dass lediglich indirekte 
Stellvertretergrößen für die Qualität der Keypoints optimiert werden~\cite{szeliski2022}.

In einer zweiten, direkteren Variante wird der Forward-Pass von XFeat so genutzt, dass dichte Karten wie 
Keypoint-Heatmaps oder Deskriptormaps in den Loss einbezogen werden können. Dadurch lässt sich die Konsistenz der 
Merkmalsdarstellung zwischen Original- und deformiertem Bild direkt bewerten. Der Generator kann so noch gezielter 
an den von XFeat genutzten Merkmalen ausgerichtet werden. Allerdings ist dieser White-Box-Ansatz deutlich 
rechenintensiver und erfordert Eingriffe in den Modellcode, da bestimmte Operationen wie die Top-K-Selektion nicht 
differenzierbar sind~\cite{detone2018superpoint,sarlin2020superglue}.

Zusammenfassend bietet Konzept~2 den Vorteil, dass der Generator durch zusätzliche strukturelle Informationen aus 
einem etablierten Feature-Detektor gesteuert wird. Dadurch kann die Robustheit der erzeugten Bilder im Hinblick auf 
nachgelagerte Matching- oder Servoing-Aufgaben verbessert werden. Nachteile sind ein erhöhter Implementierungsaufwand 
sowie das Risiko, dass bei zu starker Gewichtung des Zusatzlosses Artefakte entstehen oder der Generator Strategien 
entwickelt, die zwar die Loss-Kriterien erfüllen, aber nicht zwingend zu qualitativ hochwertigen oder realistischen 
Bildern führen. Aus diesem Grund empfiehlt sich eine moderate Gewichtung des XFeat-Losses sowie eine Kombination mit 
weiteren Verlustkomponenten wie adversarialen, rekonstruktiven oder perzeptuellen 
Losses~\cite{goodfellow2014gan,isola2018Pix2Pix}.

Nach längerer Überlegung kam die Projektgruppe zu dem Schluss, dass Konzept~2 zwar interessante Ansätze bietet, 
jedoch mit erheblichen Herausforderungen verbunden ist. Dazu zählen die Abhängigkeit von einem kombinierten Training, 
das möglicherweise nicht praktikabel umsetzbar ist, sowie die fehlende klare Definition, was im Sinne des Losses 
als „gut“ oder „schlecht“ zu bewerten ist. Hinzu kommt die Gefahr, dass die Loss-Definition im Verlauf des Trainings 
verschoben wird und der Generator sogenannte „Shortcuts“ findet, sodass der Loss reduziert werden kann, ohne dass 
tatsächlich das gewünschte Verhalten erlernt wird. Aufgrund dieser Unsicherheiten und potenziellen Fehlanreize würde 
die Projektgruppe, sollte das Thema weiterverfolgt werden, nicht primär auf Konzept~2 setzen.

\subsection{Konzept 3: Keypoint-Position}
\label{sec:Konzept3}
Das dritte Konzept verfolgt den Ansatz, den Generator anhand der Position der erkannten Features zusätzlich zu trainieren. Ziel ist es, Keypoints auf 
stabilen und strukturell relevanten Bildbereichen – vorzugsweise auf Kanten – zu belohnen, während Keypoints in reinen Textur- oder homogenen Flächen 
bestraft werden. Dazu werden die von XFeat gelieferten Keypoint-Koordinaten genutzt~\cite{potje2024xfeat}, die eine präzise Bestimmung der Lage der Features 
ermöglichen. Auf Basis dieser Koordinaten wird für jedes generierte Bild eine Kantenkarte $E(x,y)$ berechnet, beispielsweise mit differentiellen Operatoren wie Sobel, Canny oder Harris-Corner~\cite{szeliski2022}. Die Kantenstärke und -richtung an den 
Keypoint-Positionen wird anschließend ausgewertet, sodass sich ein Score ergibt, der misst, wie konsistent die Keypoints mit strukturell relevanten Kanten 
übereinstimmen.

Für die Berechnung dieses Scores können mehrere Kriterien herangezogen werden: eine ausreichend hohe Kantenstärke im Original- und im deformierten Bild, 
eine ähnliche Gradientenrichtung sowie optional die Überprüfung, ob der Punkt ein lokales Maximum entlang der Kanten-Normalen darstellt. Ein hoher Score 
bedeutet, dass der Keypoint konsistent auf einer Kante liegt, während ein niedriger Score auf eine Position in einem ungeeigneten Bereich hindeutet. Dieser 
Score wird normiert und als zusätzlicher Loss-Term in die Gesamtverlustfunktion integriert, sodass der Generator gezielt dazu angeleitet wird, Bildstrukturen 
zu erzeugen, die stabile und wiedererkennbare Features für nachgelagerte Matching- und Servoing-Aufgaben 
bereitstellen~\cite{chaumette2006visual,chaumette2007visual}.

Die Idee, Kanten als robuste Ankerpunkte für Keypoints zu nutzen, ist in der Literatur gut belegt. Klassische Verfahren wie SIFT oder SURF nutzen bereits 
Gradienteninformationen zur stabilen Feature-Detektion~\cite{lowe2004distinctive,bay2006surf}. Neuere Arbeiten wie Feature Matching Based on Consistent Edge 
Structures~\cite{mdpi2020edge} oder SiLK~\cite{gleize2023silk} zeigen ebenfalls, dass kantenbasierte Strukturen eine hohe Robustheit und Wiederholbarkeit 
für Feature-Matching bieten.

Für das geplante Einsatzszenario – die Bestimmung von Greifpositionen und Ausrichtungen an texturarmen Bauteilen – ist dieser Ansatz besonders attraktiv: 
Kanten liefern stabile geometrische Ankerpunkte, die weniger anfällig für Beleuchtungsänderungen und Oberflächenrauschen sind. Gleichzeitig bestehen 
Herausforderungen: Die Definition, wann ein Punkt „auf einer Kante“ liegt, hängt von Schwellenwerten und Kantenbreite ab, die datensatzabhängig variieren 
können. Zudem besteht das Risiko, dass der Generator künstliche Kantenartefakte erzeugt, um den Reward zu maximieren. Hier helfen Regularisierungen wie Total 
Variation oder Frequenzbestrafungen, um unnatürliche Strukturen zu verhindern. Trotz dieser Einschränkungen stellt Konzept~3 eine vielversprechende 
Möglichkeit dar, die Qualität der vom Generator erzeugten Bilder im Hinblick auf robuste und anwendungsrelevante Feature-Erkennung gezielt zu verbessern.